% ══════════════════════════════════════════════════════════════════════════════
%  lesson-macros.tex — كل ماكروهات الوثائق مشتركة (شبكة جداول بخطوط عمودية)
%  يُستدعى من common.tex — لا يحمّل أي حزمة هنا.
%
%  ملاحظة محاذاة: لا نستعمل أعمدة p{} بجوار أعمدة X — تختلف نقاط ارتكازهما
%  العمودية (بمقدار \ht\strutbox) فتنزل التسمية عن القيمة. الحل: خلايا التسمية
%  تُغلَّف في \parbox[t] (أو minipage[t]) داخل عمود c، وتتماشى مع عمود X.
% ══════════════════════════════════════════════════════════════════════════════

% ── خلايا محاذاة علوية ───────────────────────────────────────────────────────
% \celllabel{العرض}{المحتوى}   ← تسمية يمين (raggedleft)
% \cellcenter{العرض}{المحتوى}  ← وسط
\newcommand{\celllabel}[2]{\parbox[t]{#1}{\raggedleft #2}}
\newcommand{\cellcenter}[2]{\parbox[t]{#1}{\centering #2}}

% ── خلية رأس: تُستعمل مع \rowcolor{primary} (نص أبيض وسط) ──────────────────
\newcommand{\thead}[2][\linewidth]{\parbox[t]{#1}{\centering\textcolor{white}{\bfseries #2}}}

% ── خلية رأس مبسّطة (خلفية primary ونص أبيض) — للجداول الممتدة/الأفقية ──────
\newcommand{\thd}[1]{\cellcolor{primary}\color{white}\textbf{#1}}

% ── البطاقة الفنية ───────────────────────────────────────────────────────────
\newcommand{\cardsetup}{%
  \renewcommand{\arraystretch}{1.3}%
  \setlength{\arrayrulewidth}{0.4pt}%
  \arrayrulecolor{rulegray}%
}
\newcommand{\cardrow}[2]{\celllabel{3.2cm}{\textbf{\color{primary}#1}} & #2 \\ \hline}
\newcommand{\cardrowtwo}[4]{\celllabel{3.2cm}{\textbf{\color{primary}#1}} & #2 \hfill \textbf{\color{primary}#3:} #4 \\ \hline}
\newcommand{\cardtitle}[1]{%
  \begin{center}
    {\LARGE\bfseries\color{primary} #1}
  \end{center}
  \vspace{6pt}
}

% ── مراحل سير الوضعية التعلمية (مذكرة الدرس) ────────────────────────────────
\newcommand{\stageheader}{%
  \noindent\small
  \begin{tabularx}{\textwidth}{@{}|c|X|c|c|c|c|@{}}
    \hline
    \rowcolor{primary}
    \thead[1.8cm]{مرحلة الدرس} & \thead{مهمات التعلم والهدف منها} &
    \thead[2.7cm]{نشاط الأستاذ} & \thead[2.7cm]{نشاط المتعلم} &
    \thead[2.0cm]{الاستراتيجية} & \thead[1.8cm]{نوع التقويم} \\
    \hline
  \end{tabularx}\par\vspace{-0.4pt}%
}
\newcommand{\stagecontent}[6]{%
  \noindent\small
  \begin{tabularx}{\textwidth}{@{}|c|X|c|c|c|c|@{}}
    \cellcenter{1.8cm}{\textbf{\color{primary}#1}} &
    \celllabel{\linewidth}{#2} &
    \celllabel{2.7cm}{#3} &
    \celllabel{2.7cm}{#4} &
    \cellcenter{2.0cm}{\textbf{\color{secondary}#5}} &
    \cellcenter{1.8cm}{\textbf{#6}} \\
    \hline
  \end{tabularx}\par\vspace{-0.4pt}%
}

% ── مراحل سير الحصة التطبيقية (مذكرة TP) ────────────────────────────────────
\newcommand{\tpstage}[3]{%
  \noindent
  \begin{tabularx}{\textwidth}{@{}|c|c|X|@{}}
    \hline
    \rowcolor{primary}
    \thead[3.0cm]{#1} & \thead[2.0cm]{#2} & \thead{ } \\
    \hline
    \cellcenter{3.0cm}{\textbf{\color{primary}#1}} &
    \cellcenter{2.0cm}{\textbf{#2}} &
    \celllabel{\linewidth}{#3} \\
    \hline
  \end{tabularx}\par\vspace{-0.4pt}%
}

% ── جدول المهام التطبيقية ───────────────────────────────────────────────────
\newcommand{\taskheader}{%
  \noindent
  \begin{tabularx}{\textwidth}{@{}|c|X|X|@{}}
    \hline
    \rowcolor{primary}
    \thead[2.7cm]{المهمة التطبيقية} & \thead{تعليمات الأستاذ والخطوات العملية} &
    \thead{عمل المتعلم على الحاسوب} \\
    \hline
  \end{tabularx}\par\vspace{-0.4pt}%
}
\newcommand{\taskrow}[3]{%
  \noindent
  \begin{tabularx}{\textwidth}{@{}|c|X|X|@{}}
    \cellcenter{2.7cm}{\textbf{\color{primary}#1}} &
    \celllabel{\linewidth}{#2} &
    \celllabel{\linewidth}{#3} \\
    \hline
  \end{tabularx}\par\vspace{-0.4pt}%
}

% ── وثيقة التلميذ ───────────────────────────────────────────────────────────
\newcommand{\sheettitle}[4]{%
  \begin{center}
    {\Large\bfseries\color{primary} وثيقة التلميذ — الأعمال التطبيقية رقم: #1}
  \end{center}
  \vspace{2pt}
  \begin{center}
    {\large\bfseries\color{secondary} #2}
  \end{center}
  \begin{center}
    {\normalsize\color{darktext} المكان: #3 \quad$\cdot$\quad المدة الزمنية: #4}
  \end{center}
  \vspace{6pt}
}
\newcommand{\studentid}{%
  {\centering
  \begin{tabularx}{0.9\textwidth}{@{}|c|X|c|c|@{}}
    \hline
    \celllabel{2.4cm}{\textbf{\color{primary}اللقب والاسم:}} & \blank[4.8cm] &
    \celllabel{2.7cm}{\textbf{\color{primary}القسم / الفوج:}} & \blank[2.2cm] \\
    \hline
    \celllabel{2.4cm}{\textbf{\color{primary}رقم الجهاز:}} & \blank[4.8cm] &
    \celllabel{2.7cm}{\textbf{\color{primary}التاريخ:}} & \blank[2.2cm] \\
    \hline
  \end{tabularx}\par}\vspace{8pt}%
}
\newcommand{\specheader}{%
  \begingroup\centering
  \begin{tabular}{@{}|p{6.0cm}|p{9.0cm}|@{}}
    \hline
    \rowcolor{primary}
    \thead[6.0cm]{العنصر المستهدف} & \thead[9.0cm]{المواصفات المدوّنة على الجهاز} \\
    \hline
}
\newcommand{\specrow}[2]{#1 & \specdots{#2} \\ \hline}
\newcommand{\specend}{\end{tabular}\par\endgroup\vspace{8pt}}
\newcommand{\specdots}[1]{%
  \makebox[\linewidth]{\dotfill}%
  \ifnum#1>1 \\\makebox[\linewidth]{\dotfill}\fi
  \ifnum#1>2 \\\makebox[\linewidth]{\dotfill}\fi
  \ifnum#1>3 \\\makebox[\linewidth]{\dotfill}\fi
}
\newcommand{\graderow}[2]{#1 & #2 \\ \hline}

% ── ملخّص الدرس ─────────────────────────────────────────────────────────────
\newcommand{\summarytitle}[3]{%
  \begin{center}
    {\large\bfseries\color{primary} ملخّص الدرس — الأثر الكتابي على الكراس}
  \end{center}
  \vspace{2pt}
  {\cardsetup
  \centering
  \begin{tabularx}{0.9\textwidth}{@{}|c|X|@{}}
    \hline
    \celllabel{3.0cm}{\textbf{\color{primary}المجال التعلمي}} & #1 \\ \hline
    \celllabel{3.0cm}{\textbf{\color{primary}الوحدة التعلمية}} & #2 \\ \hline
    \celllabel{3.0cm}{\textbf{\color{primary}الموضوع}} & #3 \\ \hline
  \end{tabularx}\par}\vspace{8pt}
}
\newcommand{\summaryaxis}[2]{%
  \par\vspace{6pt}%
  {\normalsize\bfseries\color{primary} #1. #2}\par\vspace{3pt}%
}
\newcommand{\term}[2]{%
  \textbullet\ \textbf{\color{primary}#1:} #2\par\vspace{2pt}%
}

% ── البرنامج السنوي (قائمة، لا جدول) ────────────────────────────────────────
\newcommand{\domain}[2]{%
  \par\vspace{10pt}%
  {\large\bfseries\color{primary} المجال التعلمي #1: #2}\par\vspace{5pt}%
}
\newlist{units}{enumerate}{1}
\setlist[units,1]{%
  label=\textcolor{secondary}{\bfseries\arabic*.},
  labelwidth=1.6em, left=0.4em .. 1.8em,
  itemsep=3pt, topsep=3pt, parsep=0pt,
}
